\section{Proposed System}
\label{sec:proposed}

\subsection{System Overview}
The proposed system is a full-stack web application comprising a Python FastAPI backend and a React single-page frontend, communicating over a REST HTTP API. The backend is organized into four logical layers: (i) an ingestion and preprocessing layer that accepts an uploaded file and produces a clean, gap-filled, outlier-capped time series together with any usable exogenous features; (ii) a model orchestration layer that trains the full ensemble of candidate forecasting models concurrently and computes a composite selection score; (iii) a persistence layer that caches forecast results and conversational state to disk, keyed by a client-generated session identifier, so that concurrent users or browser tabs do not overwrite one another's results and so that state survives a backend restart; and (iv) an explanation layer that generates natural-language narration of the forecast, grounded jointly in retrieved domain knowledge and in the pipeline's own deterministic metrics. Figure~\ref{fig:architecture} depicts the overall architecture.

\subsection{Workflow Overview}
From the end user's perspective, the workflow proceeds in four steps: (1) a dataset file is uploaded and the system responds with detected columns and suggested date, target, grouping, and exogenous-feature columns; (2) the user confirms or adjusts these column selections and a forecast horizon, and optionally selects a country code for holiday-effect modeling or a grouping value to forecast an individual product, store, or category; (3) the backend executes the full training and selection pipeline and returns the leaderboard of all attempted models together with the selected best model's forecast, uncertainty interval, and, where applicable, feature importance; and (4) the user may inspect a model-comparison chart, ask the conversational assistant questions about the result, or request a structured business insights report.

\subsection{Data Pipeline}
\label{subsec:data-pipeline}
The ingestion layer accepts CSV, Microsoft Excel (\texttt{.xlsx}/\texttt{.xlsm}), and JSON files. Column detection proceeds by a combination of dtype inspection, column-name keyword matching against curated vocabularies for date-like, target-like, grouping-like, price-like, and promotion-like columns, and, for date columns specifically, an ordered cascade of parsing strategies including native pandas inference, explicit format cycling, Unix-timestamp interpretation bounded to a plausible calendar-year range, and finally mixed-format parsing as a last resort. Once a date and target column are selected, the preprocessing pipeline (i) parses and coerces the target column to a numeric type, handling currency symbols, thousands separators, and both American and European decimal conventions; (ii) sets the date column as index and aggregates duplicate timestamps by summation, consistent with the assumption that duplicate timestamps represent multiple transactions on the same period; (iii) infers the dominant sampling frequency and reindexes to a complete, gap-free calendar at that frequency; (iv) fills newly introduced and pre-existing missing values through time-aware interpolation followed by forward- and backward-filling; (v) detects and Winsorizes outliers using an interquartile-range rule; and (vi) applies light rolling-median smoothing for high-frequency (sub-daily) series. If exogenous feature columns were selected, they are aligned to the same cleaned index using an aggregation rule appropriate to their type, price-like columns by mean and promotional-flag columns by maximum, and the resulting exogenous frame is retained separately from, but index-aligned with, the cleaned target series.

\subsection{Training Pipeline}
The orchestration layer submits every registered model to a shared thread pool and trains them concurrently, subject to a per-model timeout appropriate to its statistical, machine-learning, or deep-learning category. Each model is evaluated using walk-forward, expanding-window cross-validation, with the number of folds determined adaptively from the number of available observations, subject to a configured minimum and maximum. Deep learning models are withheld entirely, rather than trained and reported unreliably, when the series has fewer observations than an empirically justified minimum for its inferred sampling frequency, and this decision is surfaced explicitly to the end user together with its justification. When exogenous features are present, a request-scoped override of the model registry substitutes exogenous-aware variants of the gradient boosting, random forest, and seasonal autoregressive models, leaving the remaining models untouched, so that the shared orchestration logic itself does not need to be aware of which individual models accept additional arguments.

\subsection{Prediction Pipeline and Model Selection}
Once every model has completed training and cross-validation, or timed out or raised an exception, the selection layer first filters to models that both completed successfully and produced a numerically usable forecast, defined as a forecast that is entirely finite and does not deviate implausibly from the historical distribution of the series. Among the surviving candidates, a composite score is computed from six components, normalized root-mean-squared error, mean absolute error, mean absolute percentage error, a discontinuity-and-volatility stability penalty, a computation-time penalty, and the coefficient of determination, described formally in Section~\ref{sec:methodology}. The model with the lowest composite score is selected, and its forecast, together with an uncertainty interval derived from its own historical cross-validation error, is returned as the pipeline's headline result, while every model's individual leaderboard entry is retained for inspection.

\subsection{Explainability Pipeline}
The explanation layer operates on the deterministic output of the prediction pipeline; it never has access to, and cannot alter, the underlying model training or selection. Two complementary mechanisms are provided. First, for the gradient boosting and random forest models, the native impurity- or gain-based feature-importance scores produced during training are exposed directly. Second, a retrieval-augmented large-language-model layer accepts a natural-language question or generates a proactive report, retrieves the most relevant passages from a curated retail-domain knowledge base using embedding similarity, and combines this retrieved context with a deterministic textual summary of the forecast's own metrics, trend classification, and, where applicable, exogenous feature usage, before generating a response. The prompt supplied to the language model explicitly enumerates the system's own computed labels and instructs the model never to re-derive or contradict them.

\subsection{User Interaction and Frontend Workflow}
The frontend is implemented as a single-page React application. It presents a sequential workflow, upload, column configuration, and results, and, once results are available, exposes tabbed views for the historical-and-forecast chart, a log-scale model-comparison chart across the full leaderboard, a forecast confidence-band chart, and, when the selected model exposes it, a feature-importance chart. A conversational assistant panel and a dedicated business insights report page are reachable from the results view. All charts are constructed from data returned directly by the backend rather than fabricated client-side, including the historical series itself and the uncertainty interval.

\subsection{API Communication}
The frontend and backend communicate exclusively through a versioned REST API implemented with FastAPI, using JSON request and response bodies. The principal endpoints are an upload endpoint that returns detected columns and suggestions; a forecast endpoint that accepts a filename, date and target column selection, forecast horizon, and optional exogenous-feature, grouping, country, and multi-target parameters, and returns the full leaderboard and selected model; a cached-result retrieval endpoint keyed by session identifier; a per-group ranking endpoint; a conversational endpoint; and a report-generation endpoint. Cross-origin resource sharing is restricted to an explicit allow-list of development origins, extensible via an environment variable for a deployed frontend origin.

\subsection{Storage and Persistence}
The system deliberately does not employ a relational or document database. Uploaded files are persisted to a dedicated directory on disk, referenced thereafter by a randomly generated filename to avoid collisions. Preprocessed series are separately persisted, both as a serialized object for programmatic reuse and as a plain comma-separated file for inspection. The retrieval-augmented explanation layer maintains a local, persistent vector store for the retail-domain knowledge-base embeddings. Finally, session-scoped state, namely the most recent forecast result and any conversational history, is held in an in-memory dictionary for low-latency access during normal operation and is additionally mirrored to a per-session file on disk, so that a backend process restart does not discard active sessions; this file-based mirror was judged sufficient for a single-process deployment and avoids introducing an external database dependency that the current deployment scope does not otherwise require.
